% BANKS-SOURCE: pdf=23 print=13 kind=problem-section id=chapter02-problems
\begin{banksproblems}

% BANKS-SOURCE: pdf=23 print=13 kind=problem id=2.1
\BanksProblem{2.1}{*}
Show that a Lorentz transformation, satisfying
\[
  \Lambda^{T}\eta\Lambda=\eta,
\]
has $\det\Lambda=\pm1$ and $|\Lambda^{0}{}_{0}|\geq 1$. The Lorentz group thus breaks up into four disconnected components depending on these two choices of sign. Only the $\Lambda^{0}{}_{0}\geq 1$, $\det\Lambda=1$ component is continuously connected to the identity. These are called \emph{proper orthochronous Lorentz transformations}. Show that we can get to the other components of the group by multiplying proper orthochronous transformations by time reversal, space reflection, or the product of these two. Show that proper orthochronous Lorentz transformations preserve the sign of the time component of time-like and null vectors in $D$ spacetime dimensions. Show that, for a space-like vector, an appropriate proper orthochronous Lorentz transformation can change the sign of the time component or set it equal to zero.

% BANKS-SOURCE: pdf=24 print=14 kind=problem id=2.2
\BanksProblem{2.2}{*}
The Dirac-picture evolution operator is defined by
\[
  U_{D}(t,t_{0})=\ee^{\ii H_{0}t}U_{S}(t,t_{0})\ee^{-\ii H_{0}t_{0}},
\]
where
\[
  \ii\partial_{t}U_{S}=(H_{0}+V)U_{S}
\]
is the equation satisfied by the Schrödinger-picture evolution operator. Show that
\[
  \ii\partial_{t}U_{D}=V(t)U_{D},
\]
where
\[
  V(t)\equiv\ee^{\ii H_{0}t}V\ee^{-\ii H_{0}t}.
\]

% BANKS-SOURCE: pdf=24 print=14 kind=problem id=2.3
\BanksProblem{2.3}{*}
Show that the solution of the Dirac-picture equation is
\[
  U_{D}(t,t_{0})=T\ee^{-\ii\int_{t_{0}}^{t}V(s)\,\dd s}.
\]

% BANKS-SOURCE: pdf=24 print=14 kind=problem id=2.4
\BanksProblem{2.4}{*}
Compute the overlap of the ground-state wave functions of a harmonic oscillator with two different frequencies. A free-bosonic field theory is just a collection of oscillators. Use your calculation to show that the overlap of the ground states for two different values of the mass is zero for any field theory, in any number of dimensions and infinite volume. Show that the overlap is zero even in finite volume if the number of space dimensions is four or greater. This is symptomatic of a more general problem. The states of two field theories, containing the same fields but with different parameters in the Lagrangian, do not live in the same Hilbert space. The formulation of field theory in terms of Green functions and functional integrals avoids this problem.

% BANKS-SOURCE: pdf=24 print=14 kind=problem id=2.5
\BanksProblem{2.5}{*}
Show that the scalar field
\[
  \phi(x)=\int\frac{\dd^{D-1}\mathbf p}{\sqrt{(2\pi)^{D-1}2\omega_{\mathbf p}}}
  \left(a(\mathbf p)\ee^{\ii p\mathbin{\cdot}x}+a^{\dagger}(\mathbf p)\ee^{-\ii p\mathbin{\cdot}x}\right)
\]
does indeed transform as a scalar, as a consequence of the transformation law of the annihilation operators. Show that it is local if and only if the particles are bosons. Compute the equal-time commutators
\[
  [\phi(t,\mathbf{x}),\phi(t,\mathbf{y})]=0,
  \qquad
  [\phi(t,\mathbf{x}),\dot\phi(t,\mathbf{y})]=\ii\delta^{D-1}(\mathbf{x}-\mathbf{y}).
\]
In the language of classical field theory, this means that the field and its time derivative are \emph{canonically conjugate}.

% BANKS-SOURCE: pdf=24 print=14 kind=problem id=2.6
\BanksProblem{2.6}{*}
Repeat the locality computation for the charged scalar field
\[
  \phi(x)=\int\frac{\dd^{D-1}\mathbf p}{\sqrt{(2\pi)^{D-1}2\omega_{\mathbf p}}}
  \left(a(\mathbf p)\ee^{\ii p\mathbin{\cdot}x}+\eta b^{\dagger}(\mathbf p)\ee^{-\ii p\mathbin{\cdot}x}\right).
\]
In this case the non-trivial computation is that of the commutator $[\phi(x),\phi^{\dagger}(y)]$. Show that the mass of the particle and that of the anti-particle (annihilated by $a$ and $b$, respectively) must be equal and that the complex number $\eta$ is a pure phase, which can be absorbed into the definition of $b$.

% BANKS-SOURCE: pdf=25 print=15 kind=problem id=2.7
\BanksProblem{2.7}{*}
Use the Dirac-picture evolution equations, with
\[
  V=-\int\dd^{D-1}\mathbf x\,\mathcal{L}_{I},
\]
to show that the interaction-picture evolution operator over infinite times is Lorentz-invariant if $\mathcal{L}_{I}$ is a local scalar operator.

% BANKS-SOURCE: pdf=25 print=15 kind=problem id=2.8
\BanksProblem{2.8}{*}
The states of a particle of mass $m$ and spin $j$ at rest transform according to the spin-$j$ representation of the rotation group.
\[
  U(R)\ket{0,k}=D^{j}_{kl}(R)\ket{0,l}.
\]
Define the states with momentum $p$ by
\[
  \ket{\mathbf p,k}=\sqrt{\frac{m}{\omega_{\mathbf p}}}\,U(L(p))\ket{0,k}.
\]
$L(p)$ is a boost that takes the rest-frame momentum into $(\omega_{\mathbf p},\mathbf{p})$. Show that the Lorentz transformation law of these states is
\[
  U(\Lambda)\ket{\mathbf p,k}=\sqrt{\frac{\omega_{\boldsymbol{\Lambda p}}}{\omega_{\mathbf p}}}\,
  D^{j}_{kl}(R_{W})\ket{\boldsymbol{\Lambda p},l},
\]
where $R_{W}=L^{-1}(\Lambda p)\Lambda L(p)$ is called the Wigner rotation. We can prove that it is a rotation by showing that it leaves the rest frame invariant. Using these transformation laws, show that one can construct fields transforming in the $[2j+1,1]$ and $[1,2j+1]$ representations of the Lorentz group (see Chapter 5)\footnote[5]{These representations are often denoted by their highest weight, rather than their dimension, and are called $[j,0]$ and $[0,j]$ representations.}, which are linear in creation and annihilation operators of these particles. Show that these fields are Hermitian and local only if we choose Fermi statistics for half-integer-spin and Bose statistics for integer-spin particles. Compute the time-ordered two-point functions for these fields and note the ultraviolet behavior of the momentum-space propagator as a function of the spin.
\setcounter{footnote}{5}

% BANKS-SOURCE: pdf=25 print=15 kind=problem id=2.9
\BanksProblem{2.9}{*}
Writing an infinitesimal Lorentz transformation as $\Lambda\simeq 1+\ii\omega^{\mu\nu}J_{\mu\nu}$, where $J_{\mu\nu}$ is a basis in the space of all $\eta$-anti-symmetric $D\times D$ matrices $(J\eta+\eta J^{T})=0$: $\eta$ is the Minkowski metric, considered as a matrix). Show that
\[
  [J_{\mu\nu},J_{\alpha\beta}]
  =\ii\left(\eta_{\mu\alpha}J_{\nu\beta}-\eta_{\nu\alpha}J_{\mu\beta}
  +\eta_{\nu\beta}J_{\mu\alpha}-\eta_{\mu\beta}J_{\nu\alpha}\right).
\]
Write this out for the individual components $J_{0i}$; in $D=4$, also write $J_{ij}=\epsilon_{ijk}J_{k}$.

% BANKS-SOURCE: pdf=25 print=15 kind=problem id=2.10
\BanksProblem{2.10}{}
Show that the stability subgroup of a null momentum\footnote[6]{Subgroup of the Poincaré group which leaves it invariant.} is isomorphic to the Euclidean group of translations and rotations in two dimensions. The only finite-dimensional representations of this group have the translation generators set to zero, while the rotation generator has a fixed eigenvalue, the helicity $h$. Following the method of induced representations we used for massive particles in the previous problem, work out the unitary representation of the Poincaré group on single massless particle states and on the corresponding creation and annihilation operators. Show that, in order to construct Lorentz-covariant local fields, we must include both signs of helicity, $\pm h$, and the helicity must be quantized in half-integer units (for fermions) and integer units (for bosons). In $D=4$, show that for helicity $\pm1$ the smallest-dimension field operator transforms like the electromagnetic field strength $F_{\mu\nu}$. Work out the analogous statement for helicity $3/2$ and $2$.
\setcounter{footnote}{6}

% BANKS-SOURCE: pdf=26 print=16 kind=problem id=2.11
\BanksProblem{2.11}{*}
Calculate the vacuum expectation value of the time-ordered exponential
\[
  \bra{0}T\ee^{\ii\int\dd^{D}x\,\phi(x)J(x)}\ket{0},
\]
for a free field. Do this by writing $\phi(x)$ as a sum of an operator involving only creation operators and another with only annihilation operators. Compute a few terms in the power-series expansion in $J$, in order to convince yourself that the answer is
\[
  \exp\!\left[
    \frac{\ii}{2}\int\dd^{D}x\,\dd^{D}y\,J(x)J(y)
    \int\frac{\dd^{D}p}{(2\pi)^{D}}
    \frac{\ee^{\ii p\mathbin{\cdot}(x-y)}}{p^{2}+m^{2}-\ii\epsilon}
  \right],
\]
which is the exponential of the second-order term. It should be easy for example to prove that all odd terms vanish. We will see a simpler derivation of this result below, in the language of functional integrals. Now consider the case
\[
  J(x)=\theta(T-t)\theta(t+T)\left[\delta^{D-1}(\mathbf{x})-\delta^{D-1}(\mathbf{x}-\mathbf{R})\right],
\]
with $T\gg R\gg1/m$, where $m$ is the mass of the field. Show that in this limit the answer has the form
\[
  \ee^{-\ii T V(R)},
\]
where the potential $V(R)$ is the Yukawa potential. We describe this result in the words \emph{static sources for a field experience a force due to exchange of virtual particles}.

\end{banksproblems}
